\documentclass[11pt]{article}

\usepackage{amsmath}
\usepackage{amssymb}
\usepackage{amsthm}
\usepackage{booktabs}
\usepackage[hmargin=1in]{geometry}
\usepackage{hyperref}

\renewenvironment{abstract}{%
  \small
  \begin{center}\bfseries\abstractname\end{center}
  \list{}{\leftmargin=3em \rightmargin=\leftmargin
          \listparindent=1.5em \itemindent=0pt
          \parsep=0pt \topsep=0pt}%
  \item\relax\ignorespaces
}{\endlist\medskip}

\newtheorem{theorem}{Theorem}
\newtheorem{lemma}[theorem]{Lemma}
\newtheorem{corollary}[theorem]{Corollary}
\newtheorem{proposition}[theorem]{Proposition}
\theoremstyle{definition}
\newtheorem{definition}[theorem]{Definition}
\theoremstyle{remark}
\newtheorem{remark}[theorem]{Remark}

\newcommand{\Rc}{\mathcal{R}}
\newcommand{\Ac}{\mathcal{A}}
\newcommand{\Pp}{P}
\newcommand{\Pm}{P^{-}}
\DeclareMathOperator{\rank}{rank}

\title{The smallest singular value of a sparse Mertens matrix}
\author{Jeffery Kline}
\date{\today}

\begin{document}
\maketitle

\begin{abstract}
Redheffer's $(0,1)$ matrix $R_n$ satisfies $\det R_n = M(n) := \sum_{k\le
n}\mu(k)$, and the Riemann hypothesis is equivalent to $M(n) =
O(n^{1/2+\varepsilon})$. A substantially sparser $(0,1)$ matrix $\Rc_n$ with the
same determinant --- about $2.61n$ nonzero entries against Redheffer's $\sim n\log
n$ --- was introduced in \cite{kline581}, and in \cite{kline584} its two dominant
\emph{eigenvalues} were shown to be $1\pm\sqrt{\pi(n)}+O(\log\log n)$. We determine
its \emph{singular values} at both ends. The largest is
$\sigma_{\max}(\Rc_n)=\sqrt{n}\,(1+O(1/n))$, forced by a single dense row, with
non-normality gap $\sigma_{\max}/|\lambda_{\max}|\asymp\sqrt{\log n}$. The smallest
satisfies
\[
  |M(n)|\cdot n^{-3/2+o(1)}\ \le\ \sigma_{\min}(\Rc_n)\ \le\
  |M(n)|\cdot n^{-4/3+o(1)},
\]
whence $\mathrm{RH}\iff\sigma_{\min}(\Rc_n)\ll_\varepsilon n^{-1+\varepsilon}$,
an equivalence that is two-sided and unconditional. The lower bound is attained,
giving $\sigma_{\min}(\Rc_n)=|M(n)|n^{-3/2+o(1)}$, as soon as
$|M(n)|/\|w\|\to0$, which RH itself supplies. The mechanism is that the entries of the vector
$\Ac^{-\mathsf T}(\mathbf 1-e_1)$, where $\Ac$ is the unit lower-triangular factor
of $\Rc_n$, are \emph{exactly} the sums $M(x,y)=\sum_{d\le x,\,\Pm(d)>y}\mu(d)$ of
the M\"obius function over integers free of small prime factors; the analytic
input is then Alladi's asymptotic $M(x,x^{1/\alpha})\sim x\rho'(\alpha)/\log y$
with $\rho$ Dickman's function \cite{alladiJNT}, which we cite rather than claim.
It gives $\|w\|=n^{1+o(1)}$ unconditionally, and that norm is what produces the
exponent $-3/2$ above.
Along the way we show $\|\Ac^{-1}\|=n^{1/2+o(1)}$ via a Neumann series that
terminates after $\sim\log n/\log\log n$ terms, hence
$\sigma_{\min}(\Ac)=n^{-1/2+o(1)}$; this answers, negatively for this family, a
question of Cheon and Kim \cite{cheonkim}. We also record why the problem is
severely pre-asymptotic: several natural conjectures suggested by computation up
to $n=10^7$ are false, and we identify the finite-range mechanism responsible.
\end{abstract}

\section{Introduction}\label{sec:intro}

Write $\mu$ for the M\"obius function and
\[
  M(n) := \sum_{k\le n}\mu(k)
\]
for the Mertens function. It is classical that the prime number theorem is
equivalent to $M(n)=o(n)$, and that the Riemann hypothesis is equivalent to
\begin{equation}\label{eq:rh}
  M(n) = O(n^{1/2+\varepsilon})\qquad\text{for every }\varepsilon>0 .
\end{equation}
Redheffer \cite{redheffer} introduced the $n\times n$ $(0,1)$ matrix $R_n$ with
$R(i,j)=1$ when $j=1$ or $i\mid j$, and proved $\det R_n = M(n)$. This places
\eqref{eq:rh} inside linear algebra, and a long line of work has studied the
spectrum of $R_n$ \cite{bfp,barrettjarvis,vaughanI,vaughanII,cardon,camargo}.

Much less is known about \emph{singular} values; Cl\'ement and Steinerberger
\cite{clementsteinerberger} say so explicitly, and treat only the largest. The
idea of routing the Riemann hypothesis through a smallest singular value appears
to originate with Bordell\`es and Cloitre \cite{bordellescloitre}, who study an
integer matrix $\Gamma_n$ with $\det\Gamma_n = n!\sum_{k\le n}\mu(k)/k$ and
deduce from an $LU$ decomposition that a lower bound
$\sigma_{\min}(U_n)\gg_\varepsilon n^{-1/2-\varepsilon}$ would imply RH; Cheon
and Kim \cite{cheonkim} give a criterion of the same shape for the triangular
factor of a Mertens equimodular matrix, and ask which matrices satisfy it. In
both cases the criterion is a \emph{sufficient condition} awaiting verification,
and in neither is the relevant singular value determined. The present paper
determines the singular values of one family at both ends, and answers the
Cheon--Kim question for it.

\subsection{The matrix}

The following much sparser $(0,1)$ matrix with the same determinant was
introduced in \cite[\S1]{kline584}. Let $\Pp(i)$ denote the largest prime factor
of $i>1$, and set
\begin{equation}\label{eq:defA}
  \Ac(i,j) :=
  \begin{cases}
    1 & \text{if } i=j,\\
    1 & \text{if } i \text{ is squarefree and } i/j=\Pp(i),\\
    0 & \text{otherwise,}
  \end{cases}
  \qquad
  \Rc_n := \Ac_n + e_1(\mathbf 1-e_1)^{\mathsf T},
\end{equation}
so that $\Rc_n$ agrees with $\Ac_n$ except that its first row is all ones. Then
$\Ac_n$ is unit lower triangular, $\det\Rc_n=M(n)$, and $\Rc_n$ has about $2.61n$
nonzero entries: one per diagonal, one per row of the all-ones row, and one for
each squarefree $i\le n$, the last contributing $6n/\pi^2+O(\sqrt n)$. By contrast
$R_n$ has $\sum_{i\le n}d(i)=n\log n+(2\gamma-1)n+O(n^{\theta})$ nonzeros. The two
dominant eigenvalues of $\Rc_n$ are
\begin{equation}\label{eq:kline1}
  r_\pm \;=\; 1\pm\sqrt{\pi(n)}+\tfrac12\,\frac{\pi_2'(n)}{\pi(n)}
  +O\!\left(\frac{\log\log^2 n}{\sqrt{n/\log n}}\right)
\end{equation}
by \cite[Thm.~1]{kline584}, where $\pi_k'(n)$ counts squarefree integers below
$n$ with exactly $k$ prime factors.

The reason $\sqrt{\pi(n)}$ appears in \eqref{eq:kline1}, rather than the $\sqrt n$
of the Redheffer matrix \cite[Thm.~1]{barrettjarvis}, is worth stating because the
same mechanism governs $\sigma_{\max}$ below. Writing $\Ac=I+S$, the matrix $S$ is
the adjacency matrix of a \emph{forest}: each squarefree $i$ has the unique parent
$i/\Pp(i)$. The children of the root $1$ are exactly the primes, so the hub degree
is $\pi(n)$; in $R_n$ the corresponding hub degree is $n$.

\subsection{Results}

Throughout, $\Pm(i)$ is the smallest prime factor of $i>1$ with
$\Pm(1)=\infty$, and
\begin{equation}\label{eq:defMxy}
  M(x,y) \;:=\; \sum_{\substack{d\le x\\ \Pm(d)>y}}\mu(d)
\end{equation}
is the M\"obius function summed over $y$-\emph{rough} integers. We write
$\alpha=\log x/\log y$ and let $\rho$ denote Dickman's function, defined by
$\rho(\alpha)=1$ on $0<\alpha\le1$ and $\rho(\alpha)=1-\int_1^\alpha
\rho(t-1)\,dt/t$ for $\alpha>1$.

Set $\mu_n := (\mu(1),\dots,\mu(n))^{\mathsf T}$ and
\begin{equation}\label{eq:defw}
  w \;:=\; \Ac_n^{-\mathsf T}(\mathbf 1 - e_1).
\end{equation}
The first of these is the first column of $\Ac_n^{-1}$; the vector $w$ collects the
\emph{other} columns, and it is the object that was missing from the analysis in
\cite{kline588}.

Our first result identifies $w$ completely.

\begin{theorem}[closed form]\label{thm:closed}
For $1\le j\le n$,
\[
  w_j =
  \begin{cases}
    1 & \text{if $j$ is not squarefree},\\
    M(n/j,\Pp(j)) & \text{if $j$ is squarefree},
  \end{cases}
\]
with the convention $\Pp(1):=1$, so that $M(n/1,\Pp(1))=M(n)$, and with the value
at $j=1$ diminished by $1$. Consequently
\[
  \|w\|^2 \;=\; \#\{j\le n:\ j\text{ not squarefree}\}
  \;+\;\bigl(M(n)-1\bigr)^2
  \;+\!\!\sum_{\substack{2\le j\le n\\ j\ \mathrm{squarefree}}}\!\! M(n/j,\Pp(j))^2 .
\]
\end{theorem}

So an inverse-matrix vector for $\Rc_n$ is, entry by entry, a rough-M\"obius sum
evaluated along the hyperbola $x=n/j$, $y=\Pp(j)$. Alladi's asymptotic for $M(x,y)$
then controls $\|w\|$:

\begin{theorem}[norm]\label{thm:norm}
$\|w\| = n^{1+o(1)}$, unconditionally.
\end{theorem}

Combining with an exact Sherman--Morrison factorisation gives the main result.

\begin{theorem}[smallest singular value]\label{thm:main}
Unconditionally,
\[
  |M(n)|\cdot n^{-3/2+o(1)}\ \le\ \sigma_{\min}(\Rc_n)\ \le\
  |M(n)|\cdot n^{-4/3+o(1)},
\]
the upper bound needing nothing beyond the prime number theorem (Remark
\ref{rem:hierarchy}). If moreover $|M(n)|/\|w\|\to0$, the two collapse to
$\sigma_{\min}(\Rc_n)=|M(n)|\cdot n^{-3/2+o(1)}$. In either case
\[
  \text{RH}\quad\Longleftrightarrow\quad
  \sigma_{\min}(\Rc_n)\ \ll_\varepsilon\ n^{-1+\varepsilon}.
\]
\end{theorem}

The equivalence is two-sided and unconditional: ``$\Longleftarrow$'' is the
unconditional lower bound rearranged, and ``$\Longrightarrow$'' may assume RH,
which supplies $|M(n)|/\|w\|\to0$ and hence the exact asymptotic. The hypothesis
$|M(n)|/\|w\|\to0$ is not known unconditionally: the best available bounds,
$\|w\|\gg_\varepsilon n^{1-\varepsilon}$ from Theorem \ref{thm:norm} and
$|M(n)|\ll n\exp(-c(\log n)^{3/5}(\log\log n)^{-1/5})$, do not cross.

We stress what is and is not
content here. That $\det\Rc_n=M(n)$ is \cite{kline581,kline584}, so
``$\mathrm{RH}\iff|\det\Rc_n|\ll n^{1/2+\varepsilon}$'' is immediate and carries
nothing new. The content of Theorem \ref{thm:main} is the conversion factor
$\|\mu_n\|\,\|w\| = n^{3/2+o(1)}$, that is, Theorem \ref{thm:norm}: without it the
equivalence is unavailable in either direction at the stated exponent.

At the other extreme, and for the unmodified factor:

\begin{theorem}[the other extreme]\label{thm:max}
$\sigma_{\max}(\Rc_n)=\sqrt n\,(1+O(1/n))$, and
$\sigma_{\max}(\Rc_n)/|\lambda_{\max}(\Rc_n)|\asymp\sqrt{\log n}$.
Moreover $\|\Ac_n^{-1}\|=n^{1/2+o(1)}$ and hence
$\sigma_{\min}(\Ac_n)=n^{-1/2+o(1)}$.
\end{theorem}

The exponent in the second assertion is sharp, but the $o(1)$ is an artefact of
the method: the terminating Neumann series bounds $\|\Ac_n^{-1}\|$ by
$\sum_r\sqrt{\pi_r'(n)}$, and that sum exceeds $\sqrt n$ by a factor
$1.55,1.62,1.67,1.72,1.76$ at $n=10^4,\dots,10^8$ --- monotone increasing, so the
argument cannot deliver $\asymp\sqrt n$ however carefully it is run. Numerically
$\|\Ac_n^{-1}\|/\sqrt n$ is flat at $0.8499$ out to $n=10^7$, and we expect
$\|\Ac_n^{-1}\|\asymp\sqrt n$; we do not prove it.

The last assertion answers a question of \cite{cheonkim} for this family; see
Remark \ref{rem:cheonkim}.

\subsection{Relation to earlier work}

In \cite{kline588} it is shown that for unit lower-triangular $L$, with $B$
obtained by deleting the first row and $l$ the first column of $L^{-1}$, one has
$\det BB^{\mathsf T}=\|l\|_2^2$; for the matrices considered here $l=\mu_n$ and
$\sqrt{\det BB^{\mathsf T}}=\sqrt{6n}/\pi+O(1)$. The same paper observes,
informally, that
\begin{quote}
  \emph{the Riemann hypothesis is an assertion that the constant $\mathbf 1^{\mathsf
  T}$ is, for all $n$, almost in the span of the rows of $B$.}
\end{quote}
That distance is $|\det \Rc_n|/\sqrt{\det BB^{\mathsf T}} = |M(n)|/\|\mu_n\| \asymp
|M(n)|n^{-1/2}$. Theorem \ref{thm:main} completes the picture quantitatively:
\begin{equation}\label{eq:compare}
  \underbrace{\operatorname{dist}(\mathbf 1^{\mathsf T},\ \mathrm{rowspan}\,B)}_{\textstyle
  =\ |M(n)|/\|\mu_n\|\ \asymp\ |M(n)|n^{-1/2}}
  \qquad\text{versus}\qquad
  \underbrace{\sigma_{\min}(\Rc_n)}_{\textstyle =\ |M(n)|/(\|\mu_n\|\|w\|)\ =\
  |M(n)|n^{-3/2+o(1)}} .
\end{equation}
The first is the distance from one particular vector to a hyperplane; the second
is the minimum over all unit directions. They differ by $\|w\|=n^{1+o(1)}$: the
matrix is a full factor of $n^{1+o(1)}$ closer to singular than that single
direction reveals, and quantifying the discrepancy is exactly what requires
Alladi's estimates.

\subsection{What is cited and what is proved}

The analytic engine is not ours. Theorem \ref{thm:alladi} below is
\cite[Thm.~1]{alladiJNT}, together with the contrast with $\Psi(x,y)$ that he
draws on p.~87 and that we record in Remark \ref{rem:buchstab};
\cite{alladiTAMS} is its companion. What is
established here is the linear algebra: Theorem \ref{thm:closed}, which identifies
$w$ with Alladi's sums, and the consequences \ref{thm:norm}--\ref{thm:max}. All
numerical claims are reproduced by the scripts in \texttt{code/}; the two
independent derivations of Theorem \ref{thm:alladi} obtained while this work was
in progress, and the prior-art audit, are recorded in \texttt{audit/reports/}.

\section{Notation and cited results}\label{sec:prelim}

For an arithmetic function $f$ we write $A(f)_n$ for the $n\times n$ matrix with
$(i,j)$ entry $f(i/j)$ when $j\mid i$ and $0$ otherwise. All matrices below are
real. We use $\|\cdot\|$ for the spectral norm and $\sigma_1\ge\cdots\ge\sigma_n$
for singular values, so $\sigma_{\min}=\sigma_n$.

Recall $\Ac=I+S$ from \eqref{eq:defA}, where
\begin{equation}\label{eq:defS}
  S(i,j)=1 \iff i \text{ squarefree},\ i\ge2,\ j=i/\Pp(i).
\end{equation}
Each row of $S$ has at most one nonzero entry, so $S$ is the adjacency matrix of a
forest in which every squarefree $i\ge2$ has the single parent $i/\Pp(i)$.

We use the following two results without proof.

\begin{theorem}[Alladi {\cite[Thm.~1]{alladiJNT}}]\label{thm:alladi}
If $y\ge x$ then $M(x,y)=1$. If $y=x^{1/\alpha}$ then
\[
  M(x,y) \;=\; \frac{x\,\rho'(\alpha)}{\log y} \;+\;\frac{y}{\log y}
  \;+\;O\!\left(\frac{x\,\alpha^2}{\log^2 y}\right),
\]
uniformly for $2\le y<x$.
\end{theorem}

Differentiating the defining relation for $\rho$ gives $\alpha\rho'(\alpha) =
-\rho(\alpha-1)$, so Theorem \ref{thm:alladi} may be written in the equivalent form
\begin{equation}\label{eq:alladi2}
  M(x,x^{1/\alpha}) = -\frac{x}{\log x}\,\rho(\alpha-1)\bigl(1+O_\alpha(1/\log y)\bigr),
\end{equation}
which is the shape used below. Note $\rho>0$ everywhere, so $M(x,y)$ is negative
and bounded away from $0$ relative to $x/\log x$ for each fixed $\alpha$.

\begin{remark}\label{rem:buchstab}
Alladi states the contrast with the smooth-number count
$\Psi(x,y)\sim x\rho(\alpha)$: the same sum weighted by $\mu$ produces Dickman's
\emph{derivative} $\rho'$ in place of $\rho$ itself \cite[p.~87]{alladiJNT}. It is
worth adding, though it is not Alladi's remark, that the unweighted rough count
$\Phi(x,y)=\#\{d\le x:\Pm(d)>y\}\sim x\,\omega(\alpha)/\log y$ involves Buchstab's
$\omega$; so the $\mu$-weighting replaces $\omega$ by $\rho'$.
\end{remark}

\begin{theorem}[{\cite[Thms.~1 and 3]{kline588}}]\label{thm:588}
Let $L$ be unit lower triangular, let $B$ be $L$ with its first row deleted, and
let $l$ be the first column of $L^{-1}$. Then $\det BB^{\mathsf T}=\|l\|_2^2$. For
$L=\Ac_n$ one has $l=\mu_n$ and
$\|\mu_n\|^2=\sum_{i\le n}|\mu(i)|=6n/\pi^2+O(\sqrt n)$.
\end{theorem}

\section{The inverse, and the vector \texorpdfstring{$w$}{w}}\label{sec:closed}

\begin{proof}[Proof of Theorem \ref{thm:closed}]
Since $S$ is nilpotent, $\Ac^{-1}=\sum_{r\ge0}(-S)^r$, and $(S^r)_{kj}$ counts
paths of length $r$ from $k$ to $j$ in the parent forest. Each vertex has at most
one parent, so there is at most one such path, and it exists precisely when $k$ is
squarefree and $j$ is obtained from $k$ by repeatedly deleting largest prime
factors. Writing $k$ squarefree with primes $p_1<\cdots<p_r$, the vertices
reachable from $k$ are exactly the initial products $p_1\cdots p_m$,
$0\le m\le r$. Hence $(\Ac^{-1})_{kj}\ne0$ if and only if $k$ is squarefree,
$j\mid k$, and every prime factor of $k/j$ exceeds $\Pp(j)$ --- which forces $j$
squarefree as well --- and in that case $(\Ac^{-1})_{kj}=(-1)^{\omega(k/j)}=\mu(k/j)$.

Now $w_j=\sum_{k\ge2}(\Ac^{-1})_{kj}$. If $j$ is not squarefree the only surviving
term is the diagonal one, giving $w_j=1$. If $j$ is squarefree, substitute
$d=k/j$: the conditions become $d\le n/j$, $d$ squarefree, and $\Pm(d)>\Pp(j)$,
whence $w_j=\sum_{d\le n/j,\,\Pm(d)>\Pp(j)}\mu(d)=M(n/j,\Pp(j))$; for $j=1$ the
constraint $k\ge2$ removes the term $d=1$.
\end{proof}

Two consequences are worth isolating, because they make part of the sum in
Theorem \ref{thm:closed} completely explicit. Write $x=n/j$ and $y=\Pp(j)$.

\begin{corollary}\label{cor:trichotomy}
For squarefree $j$: if $y\ge x$ then $w_j=1$; if $\sqrt x\le y<x$ then
$w_j = 1-(\pi(x)-\pi(y))$.
\end{corollary}

\begin{proof}
If $y\ge x$ the only $d\le x$ with $\Pm(d)>y$ is $d=1$. If $\sqrt x\le y<x$ then
such a $d$ is either $1$ or a prime in $(y,x]$, since a composite $y$-rough number
is at least $y^2\ge x$.
\end{proof}

The following two identities for $M(x,y)$ are standard and are used below; both
are immediate from unique factorisation.

\begin{lemma}\label{lem:identities}
$M(x,y)=M(x)+\sum_{p\le y}M(x/p,p)$, equivalently $M(x,y)=1-\sum_{y<p\le
x}M(x/p,p)$; and $M(x,y)=\sum_{a\le x,\ a\ y\text{-smooth}}M(x/a)$, the sum being
over \emph{all} $y$-smooth $a$, not merely squarefree ones.
\end{lemma}

\section{The norm of \texorpdfstring{$w$}{w}}\label{sec:norm}

\begin{proof}[Proof of Theorem \ref{thm:norm}]
\emph{Upper bound.} Trivially $|M(x,y)|\le x$, so by Theorem \ref{thm:closed}
$\|w\|^2\le\sum_{j\le n}(n/j)^2\le\zeta(2)n^2$, giving
$\|w\|\le (\pi/\sqrt6)\,n$.

\emph{Lower bound.} Fix $\alpha>1$ and let $j=p$ run over the primes in
$(n^{1/(\alpha+1)},2n^{1/(\alpha+1)}]$. For such $p$ we have $x=n/p$, $y=p$ and
$\log x/\log y=\alpha+O(1/\log n)$, so Theorem \ref{thm:alladi} applies at fixed
$\alpha$: its main term is $x\rho'(\alpha)/\log y\asymp_\alpha x/\log y$ and
dominates the error $O(x\alpha^2/\log^2 y)$ by a factor $\log y$. Since
$\rho>0$, $|w_p| = |M(n/p,p)| \gg_\alpha n^{\alpha/(\alpha+1)}/\log n$. By the
prime number theorem there are $\gg n^{1/(\alpha+1)}/\log n$ such primes, and by
Theorem \ref{thm:closed} the quantity $\|w\|^2$ is a sum of squares, so
\emph{any} subfamily bounds it below:
\[
  \|w\|^2 \;\gg_\alpha\; \frac{n^{1/(\alpha+1)}}{\log n}\cdot
  \frac{n^{2\alpha/(\alpha+1)}}{\log^2 n}
  \;=\;\frac{n^{(2\alpha+1)/(\alpha+1)}}{\log^3 n},
  \qquad\text{i.e.}\qquad
  \|w\| \;\gg_\alpha\; \frac{n^{1-\frac1{2\alpha+2}}}{\log^{3/2} n}.
\]
Letting $\alpha\to\infty$ gives $\|w\|\gg_\varepsilon n^{1-\varepsilon}$ for every
$\varepsilon>0$.
\end{proof}

\begin{remark}[the refined shape, open]\label{rem:shape}
Both a saddle-point analysis and the numerics of \S\ref{sec:numerics} indicate
$\|w\| = n\exp\bigl(-(1+o(1))\sqrt{\log n\log\log n}\,\bigr)$, the maximising
subfamily sitting at $\alpha\asymp\sqrt{\log n/\log\log n}$. This is \emph{not}
established. Theorem \ref{thm:alladi} does not reach that range: as Alladi notes
\cite[p.~87]{alladiJNT}, since $\rho'(\alpha)=-\exp\{-\alpha\log\alpha-\alpha\log\log\alpha+O(\alpha)\}$
for $\alpha>3$, ``the main terms of Theorem 1 are smaller than the error term when
$\alpha$ is large''. His Theorem 2 covers large $\alpha$ but only as an upper bound
with decay $\exp\{-(\alpha/2)\log\alpha\}$, half the true exponent. Determining
the constant therefore requires an improvement of the $1982$ estimates, and would
sharpen the $o(1)$ in Theorems \ref{thm:norm} and \ref{thm:main}.
\end{remark}

\section{The smallest singular value}\label{sec:sigmamin}

\begin{proposition}\label{prop:sm}
$\Rc_n^{-1}=\Ac_n^{-1}-\mu_n w^{\mathsf T}/M(n)$ whenever $M(n)\ne0$. If moreover
$|M(n)|/\|w\|\to0$, then
\begin{equation}\label{eq:factorisation}
  \sigma_{\min}(\Rc_n)\cdot\|\mu_n\|\cdot\|w\| \;=\; |M(n)|\,(1+o(1)).
\end{equation}
\end{proposition}

\begin{proof}
Sherman--Morrison applied to $\Rc_n=\Ac_n+e_1(\mathbf 1-e_1)^{\mathsf T}$ gives the
stated inverse, using $\Ac_n^{-1}e_1=\mu_n$ and $(\mathbf 1-e_1)^{\mathsf
T}\Ac_n^{-1}=w^{\mathsf T}$, together with $1+(\mathbf
1-e_1)^{\mathsf T}\Ac_n^{-1}e_1=M(n)$, which is $\det\Rc_n=M(n)$. The rank-one
term has norm $\|\mu_n\|\|w\|/|M(n)|$, while $\|\Ac_n^{-1}\|=n^{1/2+o(1)}$ by
Theorem \ref{thm:max}. Their ratio is
$(\|\mu_n\|/\|\Ac_n^{-1}\|)\cdot\|w\|/|M(n)|$, and the first factor is bounded
below by a positive constant since $\|\Ac_n^{-1}\|\ge\|\mu_n\|$; so the rank-one
term dominates precisely when $\|w\|/|M(n)|\to\infty$, which is the hypothesis.
Then $\sigma_{\min}(\Rc_n)=1/\|\Rc_n^{-1}\|$ gives \eqref{eq:factorisation}.
\end{proof}

The hypothesis is harmless in the application --- RH supplies it at once --- but it
is not known unconditionally, since the best lower bound for $\|w\|$ and the best
upper bound for $|M(n)|$ do not cross. This is the one place where the chain is
conditional, and it is why Theorem \ref{thm:main} is stated as a bracket.

\begin{proof}[Proof of Theorem \ref{thm:main}]
By Theorem \ref{thm:588}, $\|\mu_n\| = \sqrt{6n/\pi^2}\,(1+o(1)) = n^{1/2+o(1)}$,
and by Theorem \ref{thm:norm}, $\|w\|=n^{1+o(1)}$.

\emph{Lower bound.} The elementary bounds $\|w\|\le(\pi/\sqrt6)n$ and
$\|\mu_n\|\le\sqrt n$, with $\sigma_{\min}(\Rc_n)\ge|M(n)|/(\|\mu_n\|\|w\|)$ from
the rank-one representation, give $\sigma_{\min}(\Rc_n)\ge|M(n)|n^{-3/2+o(1)}$
unconditionally. \emph{Upper bound.} Remark \ref{rem:hierarchy} gives
$\sigma_{\min}(\Rc_n)\ll|M(n)|n^{-4/3+o(1)}$ from the prime number theorem alone.
If $|M(n)|/\|w\|\to0$, Proposition \ref{prop:sm} applies and substituting the two
norms into \eqref{eq:factorisation} collapses the bracket to
$\sigma_{\min}(\Rc_n)=|M(n)|n^{-3/2+o(1)}$.

For the equivalence: if RH holds then $M(n)\ll n^{1/2+\varepsilon}$, which forces
$|M(n)|/\|w\|\to0$ by Theorem \ref{thm:norm}, so the exact asymptotic is available
and $\sigma_{\min}\ll n^{-1+2\varepsilon}$. Conversely,
$\sigma_{\min}\ll_\varepsilon n^{-1+\varepsilon}$ combined with the unconditional
lower bound forces $M(n)\ll n^{1/2+\varepsilon}$, which is \eqref{eq:rh}. Note
that the reverse direction uses only the elementary upper bounds on $\|w\|$ and
$\|\mu_n\|$; the forward direction is what Theorem \ref{thm:norm} buys.
\end{proof}

\begin{remark}
$\Rc_n$ is singular exactly when $M(n)=0$, which occurs infinitely often, so
$\sigma_{\min}(\Rc_n)=0$ for infinitely many $n$ and no nontrivial unconditional
lower bound is available. Only upper bounds are meaningful, which is why test
vectors and projections are the natural tool; see Remark \ref{rem:hierarchy}.
\end{remark}

\begin{remark}[a hierarchy of cheaper bounds]\label{rem:hierarchy}
For any orthogonal projection $Q$ one has $\|\Rc_n^{-1}\|\ge\|\Rc_n^{-1}Q\|$, and
different choices of $Q$ recover successively sharper bounds at successively
higher cost. Taking the single test vector $\mu_n/\|\mu_n\|$ gives
$\sigma_{\min}\le|M(n)|/\|\mu_n\|$, which is \eqref{eq:compare} and is weaker than
the truth by exactly $\|w\|$. Taking $Q$ of rank one at an index $j=p$ with
$p\approx n^{1/3}$ --- where Corollary \ref{cor:trichotomy} applies and $w_p$ is an
\emph{exact} difference of prime counts --- already gives $\sigma_{\min}\ll
|M(n)|n^{-7/6+o(1)}$ using only the prime number theorem. Summing Corollary
\ref{cor:trichotomy} over the whole range $n^{1/3}\le p\le n^{1/2}$ gives
$\sigma_{\min}\ll|M(n)|n^{-4/3+o(1)}$, still with no input beyond PNT. Only the
last step from $n^{-4/3}$ to $n^{-3/2}$ requires Theorem \ref{thm:alladi}. In
particular, combining the third bound with the unconditional $M(n)\ll
n\exp(-c\sqrt{\log n})$ yields the clean estimate $\sigma_{\min}(\Rc_n)\ll
n^{-1/3+o(1)}$ by elementary means.
\end{remark}

\section{The largest singular value, and the unmodified factor}\label{sec:max}

\begin{proof}[Proof of Theorem \ref{thm:max}]
\emph{Largest singular value.} The first row of $\Rc_n$ has norm $\sqrt n$, so
$\sigma_{\max}\ge\sqrt n$. For the upper bound write $\Rc_n=e_1\mathbf 1^{\mathsf
T}+ (I-e_1e_1^{\mathsf T}) + S$. Each row of $S$ has at most one nonzero entry and
the column sums of $S$ are the numbers of children, the largest being $\pi(n)$ at
the root; the Schur test therefore gives $\|S\|\le\sqrt{1\cdot\pi(n)}=\sqrt{\pi(n)}$,
and the star at the root shows this is an equality. The triangle inequality gives
only $\sigma_{\max}\le\sqrt n+1+\sqrt{\pi(n)}=\sqrt n\,(1+O(1/\sqrt{\log n}))$,
which is weaker than claimed; for the stated error term argue on the Gram matrix
instead. Write
\[
  \Rc_n\Rc_n^{\mathsf T}=\begin{pmatrix} n & b^{\mathsf T}\\ b & C\end{pmatrix},
\]
the first row of $\Rc_n$ being $\mathbf 1^{\mathsf T}$. For $k\ge2$, $b_k$ is the
number of nonzero entries in row $k$, which is $2$ if $k$ is squarefree and $1$
otherwise; hence
\[
  \|b\|^2=4\bigl(Q(n)-1\bigr)+\bigl(n-Q(n)\bigr)=3Q(n)+n-4\ \le\ 4n,
\]
where $Q(n)=\#\{k\le n:\ k\text{ squarefree}\}$. Also
$C=BB^{\mathsf T}$ with $B$ the last $n-1$ rows of $\Rc_n$, so
$\|C\|\le(1+\sqrt{\pi(n)})^2=o(n)$. If $\lambda=\lambda_{\max}(\Rc_n\Rc_n^{\mathsf
T})\ge n$ then $\lambda>\|C\|$ for large $n$, and the Schur complement identity
for the top-left entry gives
\[
  \lambda = n + b^{\mathsf T}(\lambda I-C)^{-1}b
  \ \le\ n+\frac{\|b\|^2}{\lambda-\|C\|}\ =\ n+O(1).
\]
Hence $\sigma_{\max}=\sqrt{n+O(1)}=\sqrt n\,(1+O(1/n))$. (Numerically
$\lambda_{\max}(\Rc_n\Rc_n^{\mathsf T})-n\to 1+18/\pi^2=2.8238$, which is
$\lim\|b\|^2/n$, so the $O(1)$ is attained and not merely an upper bound.)

\emph{Non-normality.} By \eqref{eq:kline1}, $|\lambda_{\max}(\Rc_n)| =
(1+o(1))\sqrt{\pi(n)}$, and $\sqrt{n/\pi(n)}\asymp\sqrt{\log n}$.

\emph{The unmodified factor.} Let $\omega_{\max}(n)$ be the largest number of prime
factors of a squarefree integer at most $n$, so $\omega_{\max}(n)\sim\log
n/\log\log n$. Each application of $S$ deletes one prime factor, hence
$S^{\,\omega_{\max}(n)+1}=0$ and the Neumann series
$\Ac_n^{-1}=\sum_{r\le\omega_{\max}(n)}(-S)^r$ \emph{terminates}. Row sums of $S^r$
are at most $1$ and its column sums are $\pi_r'(n)$, the number of squarefree $k\le
n$ with $\omega(k)=r$; so $\|S^r\|\le\sqrt{\pi_r'(n)}$ and
$\|\Ac_n^{-1}\|\le\sum_r\sqrt{\pi_r'(n)} = n^{1/2+o(1)}$. Conversely
$\|\Ac_n^{-1}\|\ge\|\Ac_n^{-1}e_1\|=\|\mu_n\|\asymp\sqrt n$ by Theorem
\ref{thm:588}. Finally $\sigma_{\min}(\Ac_n)=1/\|\Ac_n^{-1}\|$.
\end{proof}

\begin{remark}[Bordell\`es--Cloitre]\label{rem:bc}
The strategy of bounding a smallest singular value from below in order to reach
PNT or RH is due to \cite{bordellescloitre}. There $U_n$ is the upper triangular
factor of $\Gamma_n$, with $U_n^{-1}=(\theta_{ij})$ given explicitly in terms of
iterated M\"obius functions, and Corollary 2.7 of that paper states
$\bigl|\sum_{k=i}^n \mu_i(k)/k\bigr|\le 1/(n\sigma_n)$, so that
$\sigma_n\gg_\varepsilon n^{-1+\varepsilon}$ suffices for PNT and
$\sigma_n\gg_\varepsilon n^{-1/2-\varepsilon}$ suffices for RH. The authors
observe that the general bounds available for smallest singular values of
triangular matrices --- for instance $\sigma_n\ge\min_i|a_{ii}|/2^{\,n-1}$ under a
diagonal dominance hypothesis --- are exponentially weak and ``still very far from
the PNT''.

Theorem \ref{thm:max} is a case where the arithmetic structure does determine the
order exactly: for the triangular factor $\Ac_n$ considered here,
$\sigma_{\min}(\Ac_n)\asymp n^{-1/2}$, obtained not from a general matrix
inequality but from the fact that the Neumann series for $\Ac_n^{-1}$ terminates
after $\sim\log n/\log\log n$ terms. We note in passing that this is precisely
the exponent appearing in the RH threshold of \cite{bordellescloitre}, though for
a different matrix, so the two do not combine.
\end{remark}

\begin{remark}[a question of Cheon and Kim]\label{rem:cheonkim}
For a Mertens equimodular matrix $L_n+uv^{\mathsf T}$, since $\det L_n=1$ and
$uv^{\mathsf T}$ has singular values $\sqrt{n-1},0,\dots,0$, the determinantal
inequality of Li and Mathias \cite{limathias} gives $|M(n)|\le
1+\sqrt{n-1}/\sigma_{\min}(L_n)$; Cheon and Kim \cite[Thm.~4.2]{cheonkim} deduce
that $1+\sqrt{n-1}/\sigma_{\min}(L_n)=O(n^{1/2+\varepsilon})$ implies the Riemann
hypothesis, and ask which matrices satisfy this. Theorem \ref{thm:max} answers the
question negatively for the present family: with
$\sigma_{\min}(\Ac_n)=n^{-1/2+o(1)}$ the quantity is $n^{1+o(1)}$, which is not
$O(n^{1/2+\varepsilon})$, so the resulting bound $|M(n)|\le n^{1+o(1)}$ is trivial. Note their criterion concerns the
\emph{unmodified} factor and is one-directional, whereas Theorem \ref{thm:main}
concerns the modified matrix and is an equivalence.
\end{remark}

\section{Numerical remarks}\label{sec:numerics}

All values below are produced by the scripts in \texttt{code/}, using only NumPy
(\texttt{checkB.py} additionally uses \texttt{mpmath}): the $O(n\log n)$
quantities by \texttt{wnorm.py} and \texttt{massprofile.py}, and every
singular-value, eigenvalue and nilpotency figure by \texttt{spectra.py}, which
builds the matrices explicitly and calls a dense SVD. Two independent
implementations were used throughout: a dense construction with a direct linear
solve, and an $O(n\log n)$ evaluator built from Theorem \ref{thm:closed} and
Corollary \ref{cor:trichotomy}. They agree to residual exactly $0$ at
$n=500,1000,2000$, which is the check on Theorem \ref{thm:closed}.

\begin{table}[h]\centering
\begin{tabular}{rrrr}
\toprule
$n$ & $\|w\|$ & $\|w\|/n$ & local exponent \\
\midrule
$10^5$ & $2289.3$   & $0.022893$ & --- \\
$3\cdot10^5$ & $5211.8$   & $0.017373$ & $0.7488$ \\
$10^6$ & $12850.5$  & $0.012851$ & $0.7496$ \\
$3\cdot10^6$ & $29519.5$  & $0.009840$ & $0.7570$ \\
$10^7$ & $73972.3$  & $0.007397$ & $0.7630$ \\
$3\cdot10^7$ & $172026.3$ & $0.005734$ & $0.7683$ \\
\bottomrule
\end{tabular}
\caption{The norm $\|w\|$. The local exponent climbs steadily; Theorem
\ref{thm:norm} asserts it tends to $1$.}\label{tab:w}
\end{table}

Verification of \eqref{eq:factorisation}, computed from the dense matrix: the
ratio $\sigma_{\min}(\Rc_n)\|\mu_n\|\|w\|/|M(n)|$ equals $0.9896$, $1.0038$,
$1.0047$ at $n=500,1000,2000$. Over the grid $n=200,225,\dots,3200$ (excluding the
$n$ with $M(n)=0$, where $\Rc_n$ is singular) it oscillates within
$[0.987,1.027]$, the extremes occurring at $n=300$ and $n=225$; the deviation is
bounded and shows no trend, consistent with the $1+o(1)$ of Proposition
\ref{prop:sm}, but it is not monotone and does not visibly contract over this
range. For $\|\mu_n\|$, the ratio $\|\mu_n\|/\sqrt n$ equals $0.77971$ at $n=10^5$
and $0.77970$ at $n=3\cdot10^5,10^6,3\cdot10^6,10^7$, against
$\sqrt{6/\pi^2}=0.779697$: the squarefree count is already at its asymptotic
density to five figures across this range.

For Theorem \ref{thm:max}: $\sigma_{\max}(\Rc_n)/\sqrt n$ equals
$1.00387,1.00191,1.00094,1.00047,1.00023$ at $n=400,\dots,6400$, the excess
halving at each doubling; $\sigma_{\max}/|\lambda_{\max}|$ equals
$1.906,2.087,2.262,2.440,2.592$ against $\sqrt{\log n}=2.448,\dots,2.960$; the
nilpotency index of $S$ is $5,5,5,6$ at $n=500,\dots,3000$;
$\|\Ac_n^{-1}\|/\sqrt n$ equals $0.853,0.851,0.849,0.850$; and
$\sigma_{\min}(\Ac_n)\sqrt n$ equals $1.1717,1.1753,1.1773,1.1760$. The quantity
of Remark \ref{rem:cheonkim} equals $427,851,1699,2552$ against the required
$O(n^{1/2+\varepsilon})=22,32,45,55$.

For scale, the same non-normality ratio computed for Redheffer's own $R_n$ (built
directly, and checked against $\det R_n=M(n)$) equals
$1.2243,1.2475,1.2686,1.2875$ at $n=400,800,1600,3200$, against
$1.9063,2.0871,2.2623,2.4395$ for $\Rc_n$. Both are far below $\sqrt n$, and at
these sizes the sparser matrix is the more non-normal of the two; we draw no
asymptotic conclusion from four points.

\subsection{The problem is severely pre-asymptotic}\label{sec:preasymptotic}

We record this because computation up to $n=10^7$ is actively misleading here, and
several natural conjectures it suggests are false.

Table \ref{tab:w} makes $\|w\|$ look like a clean power of $n$: over
$10^5\le n\le10^6$ the ratio $\|w\|/n^{3/4}$ is stable to four significant figures
($0.40711,0.40658,0.40637$). It is not a power. The ratio breaks at $3\cdot10^6$
($0.40951$) and the local exponent continues to climb past $3/4$, consistent with
Theorem \ref{thm:norm}. Likewise, the normalisation $\|w\|(\log
n)^{1.13}/n^{5/6}$ appears stable to under $1\%$ per decade while in fact
increasing monotonically at every step.

The mechanism is a \emph{corner}. Writing $\beta=\log j/\log n$, the mass of
$\|w\|^2$ concentrates at $\beta=1/3$, equivalently $\alpha=2$ --- exactly the
threshold of Corollary \ref{cor:trichotomy}, below which the $y$-rough numbers
under $x$ are still only $1$ and the primes. Dickman's $\rho$ has a corner at $1$,
so the relevant integrand has a jump in its derivative there, and the maximiser is
\emph{pinned} at $\beta=1/3$ until $\log n$ exceeds an explicit threshold. The
measured mass-mean $\beta$ is $0.3368,0.3364,0.3361,0.3351$ at
$n=10^5,\dots,10^7$: already drifting, and accelerating. The crossover at which the
$\alpha=3$ subfamily overtakes the $\alpha=2$ subfamily lies near $n=10^{12}$, far
beyond any feasible computation.

The practical consequence is that the exponent in Theorem \ref{thm:norm} cannot be
inferred numerically, and any conjecture about $\|w\|$ or $\sigma_{\min}(\Rc_n)$
based on data below $10^8$ should be treated with suspicion.

\section{Open problems}\label{sec:open}

\begin{enumerate}
\item \textbf{The refined shape.} Establish $\|w\| = n\exp(-(1+o(1))\sqrt{\log
n\log\log n})$, and thereby replace the $o(1)$ in Theorem \ref{thm:main} by an
explicit function. By Remark \ref{rem:shape} this requires extending Theorem
\ref{thm:alladi} to $\alpha\asymp\sqrt{\log x/\log\log x}$ --- the $\mu$-weighted
analogue of Hildebrand's range for $\Psi(x,y)\sim x\rho(\alpha)$ \cite{hildtenen}
--- or improving the decay in \cite[Thm.~2]{alladiJNT} from
$\exp\{-(\alpha/2)\log\alpha\}$ to $\exp\{-\alpha\log\alpha\}$.
\item \textbf{Other families.} By Remark \ref{rem:cheonkim} the present family
fails the Cheon--Kim criterion because its triangular factor is poorly
conditioned, $\sigma_{\min}(\Ac_n)\asymp n^{-1/2}$. Is there a family in which
$\sigma_{\min}(L_n)\gg n^{-\varepsilon}$? Theorem \ref{thm:max} shows the
obstruction here is that $\|\Ac_n^{-1}\|$ is already attained by its first column,
namely $\|\mu_n\|$; any candidate family must avoid that. The same question for
the matrix $U_n$ of \cite{bordellescloitre} appears to be open, and Theorem
\ref{thm:max} suggests that a termination argument for the relevant Neumann
series, rather than a general triangular-matrix inequality, is the tool to try.
\item \textbf{Deeper singular values.} \eqref{eq:kline1} and Theorem
\ref{thm:max} describe the extremes. The forest structure suggests that successive
shells --- squarefree integers with $2$, then $3$ prime factors --- govern the next
singular values, with $\pi_2'(n)$ already visible in \eqref{eq:kline1}.
\end{enumerate}

\begin{thebibliography}{99}

\bibitem{alladiJNT}
K.~Alladi, \emph{Asymptotic estimates of sums involving the Moebius function},
J.\ Number Theory 14 (1982) 86--98.

\bibitem{alladiTAMS}
K.~Alladi, \emph{Asymptotic estimates of sums involving the Moebius function.
II}, Trans.\ Amer.\ Math.\ Soc.\ 272 (1982) 87--105.

\bibitem{bordellescloitre}
O.~Bordell\`es, B.~Cloitre, \emph{A matrix inequality for M\"obius functions},
J.\ Inequal.\ Pure Appl.\ Math.\ 10 (2009), art.\ 62, 17 pp.

\bibitem{bfp}
W.~W.~Barrett, R.~W.~Forcade, A.~D.~Pollington, \emph{On the spectral radius of a
$(0,1)$ matrix related to Mertens' function}, Linear Algebra Appl.\ 107 (1988)
151--159.

\bibitem{barrettjarvis}
W.~W.~Barrett, T.~J.~Jarvis, \emph{Spectral properties of a matrix of Redheffer},
Linear Algebra Appl.\ 162 (1992) 673--683.

\bibitem{camargo}
A.~Pierro de Camargo, \emph{Dirichlet matrices: determinants, permanents and the
Factorisatio Numerorum problem}, Linear Algebra Appl.\ 628 (2021) 115--129.

\bibitem{cardon}
D.~A.~Cardon, \emph{Matrices related to Dirichlet series}, J.\ Number Theory 130
(2010) 27--39.

\bibitem{cheonkim}
G.-S.~Cheon, H.~Kim, \emph{Mertens equimodular matrices of Redheffer type},
Linear Algebra Appl.\ 572 (2019) 252--272.

\bibitem{clementsteinerberger}
F.~Cl\'ement, S.~Steinerberger, \emph{On the largest singular vector of the
Redheffer matrix}, Linear Algebra Appl.\ 725 (2025) 96--114.

\bibitem{debruijn}
N.~G.~de Bruijn, \emph{On the number of positive integers $\le x$ and free of
prime factors $>y$}, Nederl.\ Akad.\ Wetensch.\ Proc.\ Ser.\ A 54 (1951) 50--60.

\bibitem{dickman}
K.~Dickman, \emph{On the frequency of numbers containing prime factors of a
certain relative magnitude}, Ark.\ Mat.\ Astr.\ Fys.\ 22 (1930) 1--14.

\bibitem{hildtenen}
A.~Hildebrand, G.~Tenenbaum, \emph{Integers without large prime factors},
J.\ Th\'eor.\ Nombres Bordeaux 5 (1993) 411--484.

\bibitem{kline581}
J.~Kline, \emph{A sparser matrix representation of the Mertens function}, Linear
Algebra Appl.\ 581 (2019) 354--366.

\bibitem{kline584}
J.~Kline, \emph{On the eigenstructure of sparse matrices related to the prime
number theorem}, Linear Algebra Appl.\ 584 (2020) 409--430.

\bibitem{kline588}
J.~Kline, \emph{Bordered Hermitian matrices and sums of the M\"obius function},
Linear Algebra Appl.\ 588 (2020) 224--237.

\bibitem{limathias}
C.-K.~Li, R.~Mathias, \emph{The determinant of the sum of two matrices},
Bull.\ Austral.\ Math.\ Soc.\ 52 (1995) 425--429.

\bibitem{redheffer}
R.~Redheffer, \emph{Eine explizit l\"osbare Optimierungsaufgabe}, Numerische
Methoden bei Optimierungsaufgaben, Band 3, Birkh\"auser, 1977, pp.\ 213--216.

\bibitem{vaughanI}
R.~C.~Vaughan, \emph{On the eigenvalues of Redheffer's matrix I}, in: Number
Theory with an Emphasis on the Markoff Spectrum, Dekker, 1993, pp.\ 283--296.

\bibitem{vaughanII}
R.~C.~Vaughan, \emph{On the eigenvalues of Redheffer's matrix II}, J.\ Austral.\
Math.\ Soc.\ 60 (1996) 260--273.

\end{thebibliography}

\end{document}
